\section{Related Work}
\label{sec:related}
The relevant distinctions concern the partition rule, when it changes, and the memory that holds the propagated operator.

\noindent\textbf{Conserved Operator Subspaces.}
Symmetries and invariant subspaces are established tools in quantum simulation.
Bravyi et al. identify binary symmetries of fermionic Hamiltonians and remove qubits after selecting a symmetry sector~\cite{bravyi2017tapering}.
West et al. describe a commutator graph whose connected components identify invariant operator subspaces under a set of generators~\cite{west2025graph}.
These results establish the algebraic basis for independent evolution and preclude a novelty claim based on parity conservation alone.
\SectorQ uses binary rows as labels on Pauli strings, with the ordinary binary inner product expressing invariance under generator XOR.
The labels are not symmetry eigenvalues of a quantum state, and their sectors need not be individual connected components of the commutator graph.
The system derives sufficiently fine labels for a finite window and replaces them when the next window has different generators.
This construction supplies an executable partition for distributed merge and truncation without requiring a useful nontrivial partition conserved by the entire circuit.

\noindent\textbf{Fixed Linear Partitions.}
The closest implementation-level comparison is paulistrings-rs, whose partition label is a binary matrix times the Pauli bit vector~\cite{paulistrings}.
Its linear map identifies one communication partner for each nonlocal generator delta, and its local coset execution already exploits independent groups of strings.
The documented partition rows remain fixed during a run, and automatic row selection that accounts for both balance and remote work is listed as an open problem.
The parallel iQCC implementation also uses a bitmask to assign Pauli words to processes and computes the destination of a product through XOR~\cite{genin2025chemistry}.
A predictable destination can still be remote.
\SectorQ requires every chosen row to annihilate every generator in the current window, placing all resulting merge partners in one sector throughout that window.
It then selects rows again for the next window and schedules sectors from host memory.
Our same-engine fixed-row ablation isolates this temporal choice, while the paulistrings-rs comparison measures complete implementations with matched truncation.

\noindent\textbf{Distributed CPU Propagation.}
ORQA demonstrates large-scale propagation on Fugaku with up to one trillion Pauli strings and $2^{17}$ processes~\cite{orqa}.
Its distribution map is a structured enumeration of string bits designed to limit communication partners while maintaining balance.
It is not a generic hash, and the sparsity of its communication is central to its scaling.
The map still distributes updates as the propagation advances through generators, rather than deriving a partition conserved by a window.
MonoProp provides distributed Pauli and Majorana propagation through a C++ implementation with Python interfaces~\cite{monoprop}.
These systems establish that large operators can reside in distributed CPU memory.
\SectorQ addresses the additional requirement of executing sectors through GPU kernels while the operator remains in host memory.
Our evaluation labels the ORQA implementation as a reproduction, identifies the CPU resources used, and treats MonoProp's different truncation rule separately from comparisons with identical coefficient decisions.

\noindent\textbf{Propagation Libraries and GPU Kernels.}
Sparse Pauli dynamics and PauliPropagation.jl develop efficient observable simulation and truncation methods~\cite{begusic2024fast,rudolph2025pauli}.
Quantinuum's SPD provides GPU execution, and Qiskit's Pauli propagation package supports Clifford operations, Pauli rotations, and noise channels~\cite{spd,qiskitpauliprop}.
These libraries improve the cost and coverage of propagation within their supported execution environments.
NVIDIA's \cuPP supplies GPU operations on sparse Pauli expansions and a distributed execution mode~\cite{cupauliprop}.
\SectorQ uses its single-GPU operations for arithmetic and compares against its original-circuit distributed mode.
Consequently, the contribution concerns the locality and storage of those operations, not a replacement arithmetic kernel.
The measured element sequence separates Clifford preprocessing and deferred deduplication from window-sector execution.
The single-GPU comparison also retains cases in which another implementation is faster.
This distinction identifies the system benefit as the ability to execute growing operators across nodes without requiring the complete local expansion to remain in GPU memory.

\noindent\textbf{Repeated Evaluations and Gradients.}
PADO-Pauli compiles circuits and observables into sparse matrices that are reused for batches of parameter values and supports gradients through PyTorch~\cite{kim2026pado}.
Its reuse of a compiled representation targets repeated evaluations during variational training.
Backpropagating Pauli Propagation develops gradient evaluation whose computational complexity matches forward propagation and reduces its memory requirements through reversibility~\cite{lin2026backpropagating}.
These objectives differ from the single forward propagation measured in our capacity and scaling experiments.
\SectorQ follows the surviving terms of one specified circuit and cutoff, with new rows chosen from those live terms at window boundaries.
It neither assumes a fixed sparse program shared by parameter batches nor evaluates derivatives in the reported experiments.
The distinction is relevant because a faster batched or gradient evaluation is not directly comparable to the time of a single forward run with the same coefficient truncation.
The methods address complementary workloads, and their integration would require a separate numerical and performance evaluation.

\noindent\textbf{Staged State Vector Simulation.}
State vector simulators already reduce communication through groups of gates.
H\"aner and Steiger schedule gates to reduce communication in a distributed simulation of 45 qubits~\cite{haner2017petabyte}.
Atlas partitions circuits into stages with an integer program, selects kernels with dynamic programming, and supports host-memory offloading~\cite{xu2024atlas}.
Stage boundaries and host storage are therefore not new independently.
Their state vector has a fixed index space, whereas Pauli propagation creates strings, combines equal strings, and removes coefficients during execution.
A useful partition must preserve these merges before truncation, and the sizes of its parts depend on the surviving operator.
\SectorQ derives this locality from the window generators and obtains execution granularity from a sample of the live terms.
An oversized sector can be subdivided only by rows valid for its remaining operations.
This constraint connects host-memory execution to the algebraic partition and distinguishes the proposed mechanism from applying a state-vector stage schedule to a sparse operator.
